\documentclass{article}
\usepackage{amsmath, amssymb, amsthm}
\usepackage{hyperref}
\usepackage{geometry}
\geometry{margin=1in}

\title{Erd\H{o}s Problem \#184}
\date{Last updated: 2025-08-31}
\author{Source: erdosproblems.com}

\begin{document}
\maketitle

\noindent\textbf{Status:} OPEN \quad \textbf{No prize}

\noindent\textbf{Tags:} graph theory, cycles

\medskip

\noindent\textbf{Problem Statement:}

Any graph on $n$ vertices can be decomposed into $O(n)$ many edge-disjoint cycles and edges.




Conjectured by Erd\H{o}s and Gallai, who proved that $O(n\log n)$ many cycles and edges suffices (see Section 5 of \cite{EGP66}). The graph $K_{3,n-3}$ shows that at least $(1+c)n$ many cycles and edges are required, for some constant $c&gt;0$. In \cite{Er71} Erd\H{o}s suggests that only $n-1$ many cycles and edges are required if we do not require them to be edge-disjoint. The best bound available is due to Buci\'{c} and Montgomery \cite{BM22}, who prove that $O(n\log^*n)$ many cycles and edges suffice, where $\log^*$ is the iterated logarithm function. Conlon, Fox, and Sudakov \cite{CFS14} proved that $O_\epsilon(n)$ cycles and edges suffice if $G$ has minimum degree at least $\epsilon n$, for any $\epsilon&gt;0$. See also [583] for an analogous problem decomposing into paths, and [1017] for decomposing into complete graphs.




References

[BM22] Buci\'C, M. and Montgomery, R., Towards the Erd\H{o}s-Gallai Cycle Decomposition Conjecture . arXiv:2211.07689 (2022).

[CFS14] Conlon, David and Fox, Jacob and Sudakov, Benny, Cycle packing . Random Structures Algorithms (2014), 608-626.

[EGP66] Erd\H{o}s, Paul and Goodman, A. W. and P\'{o}sa, Lajos, The representation of a graph by set intersections . Canadian J. Math. (1966), 106-112.

[Er71] Erd\H{o}s, P., Some unsolved problems in graph theory and combinatorial analysis . Combinatorial Mathematics and its Applications (Proc. Conf., Oxford, 1969) (1971), 97-109.

\noindent\textbf{Related OEIS sequences:} possible


\bigskip
\noindent\small{Source: \url{https://www.erdosproblems.com/184}}
\end{document}
